% GENERATED FILE. Edit canonical lessons, then run npm run generate:books.
% canonical-source-hash: fnv1a64:4abe3e7ed1433715
% canonical-lessons: ZH-C214-wenti, ZH-C214-gaosu, ZH-C214-mingbai, ZH-C214-banfa, ZH-C214-zenmeban

\chapter{Asking for Help}
\label{ch:zh-asking-for-help}
\hlchaptermodality{214}

\begin{chapteropening}
\textbf{By the end of this chapter:} \emph{I can say I have a question, ask someone to tell me, say I understand, and ask what to do.}

Complete the last lesson of chapter 214: I can say I have a question, ask someone to tell me, say I understand, and ask what to do.
\end{chapteropening}

\section[wèntí]{wèntí — question; problem}

\label{lesson:ZH-C214-wenti}



\begin{quote}
Before the new one: say the Mandarin for admission ticket, then the Mandarin for have a safe journey.
\end{quote}

\subsection*{You'll want to know: wèntí}
\emph{wèntí} — \textquotedblleft{}question; problem\textquotedblright{}.

\begin{quote}
\emph{wǒ yǒu yí ge wèntí} — I've got a question.
\end{quote}

Its characters come later; for now it is a word for the ear.

\begin{cousinweb}
\emph{wèntí} is built from \zh{问} 'to ask' and \emph{tí} 'a heading', both inherited from Middle Chinese, \emph{mjunH} and \emph{dej} (Guangyun, in Baxter's transcription), which glosses the second 'the forehead' and 'the title of a book'.
\end{cousinweb}

\begin{grammarlens}[title={What you've built}]
One more word for saying no and asking.
\end{grammarlens}

\subsection*{Your turn}
\begin{itemize}
\raggedright
  \item \emph{Say it:} \emph{wèntí}
  \item \emph{Say it:} \emph{wèntí}, once more
  \item \emph{Say it:} \emph{wǒ yǒu yí ge wèntí} — I've got a question.
  \item \emph{Recall:} say the Mandarin for admission ticket, then the Mandarin for have a safe journey, then say \emph{wèntí} again
\end{itemize}

\begin{tcolorbox}[breakable,colback=teal!4,colframe=teal!35!black,title={Before you move on}]
What does wèntí mean? (\textquotedblleft{}Question; problem\textquotedblright{}.) Say it once more.
\end{tcolorbox}
\section[gàosu]{gàosu — to tell}

\label{lesson:ZH-C214-gaosu}



\begin{quote}
Before the new one: say the Mandarin for have a safe journey, then the Mandarin for problem.
\end{quote}

\subsection*{You'll want to know: gàosu}
\emph{gàosu} — \textquotedblleft{}to tell\textquotedblright{}.

\begin{quote}
\emph{qǐng gàosu wǒ nǐ de hàomǎ} — Please tell me your number.
\end{quote}

Its characters come later; for now it is a word for the ear.

\begin{cousinweb}
\emph{gàosu} is built from \emph{gào} 'to report' and \emph{su} 'to tell', both inherited from Middle Chinese, \emph{kawH} and \emph{suH} (Guangyun, in Baxter's transcription); in the word the second loses its tone.
\end{cousinweb}

\begin{grammarlens}[title={What you've built}]
One more word for saying no and asking.
\end{grammarlens}

\subsection*{Your turn}
\begin{itemize}
\raggedright
  \item \emph{Say it:} \emph{gàosu}
  \item \emph{Say it:} \emph{gàosu}, once more
  \item \emph{Say it:} \emph{qǐng gàosu wǒ nǐ de hàomǎ} — Please tell me your number.
  \item \emph{Recall:} say the Mandarin for have a safe journey, then the Mandarin for problem, then say \emph{gàosu} again
\end{itemize}

\begin{tcolorbox}[breakable,colback=teal!4,colframe=teal!35!black,title={Before you move on}]
What does gàosu mean? (\textquotedblleft{}To tell\textquotedblright{}.) Say it once more.
\end{tcolorbox}
\section[\texorpdfstring{\zh{明白} (míngbai)}{明白 (míngbai)}]{\zh{明白} (míngbai) — to grasp a meaning; to understand; clear}

\label{lesson:ZH-C214-mingbai}



\begin{quote}
Before the new one: say the Mandarin for problem, then the Mandarin for to tell someone.
\end{quote}

\subsection*{You'll want to know: \zh{明白}}
\textbf{\zh{明白}} — \emph{míngbai} — \textquotedblleft{}to grasp a meaning; to understand; clear\textquotedblright{}.

\begin{quote}
\textbf{\zh{我明白你很关心我}} \emph{wǒ míngbai nǐ hěn guānxīn wǒ} — I know how much you care about me.
\end{quote}

\begin{cousinweb}
\textbf{\zh{明白}} (\emph{míngbai}) is built from \zh{明} 'bright' and \zh{白} 'white', both inherited from Middle Chinese, \emph{mjaeng} and \emph{baek} (Guangyun, in Baxter's transcription): bright and clear. In the word the second loses its tone.
\end{cousinweb}

\begin{grammarlens}[title={What you've built}]
One more word for saying no and asking.
\end{grammarlens}

\subsection*{Your turn}
\begin{itemize}
\raggedright
  \item \emph{Say it:} \emph{míngbai}
  \item \emph{Say it:} \emph{míngbai}, once more
  \item \emph{Say it:} \emph{wǒ míngbai nǐ hěn guānxīn wǒ} — I know how much you care about me.
  \item \emph{Recall:} say the Mandarin for problem, then the Mandarin for to tell someone, then say \emph{míngbai} again
\end{itemize}

\begin{tcolorbox}[breakable,colback=teal!4,colframe=teal!35!black,title={Before you move on}]
What does \zh{明白} mean? (\textquotedblleft{}To grasp a meaning; to understand; clear\textquotedblright{}.) Say it once more.
\end{tcolorbox}
\section[bànfǎ]{bànfǎ — a way to deal with it}

\label{lesson:ZH-C214-banfa}



\begin{quote}
Before the new one: say the Mandarin for to tell someone, then the Mandarin for to grasp a meaning.
\end{quote}

\subsection*{You'll want to know: bànfǎ}
\emph{bànfǎ} — \textquotedblleft{}a way to deal with it\textquotedblright{}.

\begin{quote}
\emph{wǒ méiyǒu bànfǎ} — There's nothing I can do.
\end{quote}

Its characters come later; for now it is a word for the ear.

\begin{cousinweb}
\emph{bànfǎ} is built from \emph{bàn} 'to handle' and \emph{fǎ} 'a method', both inherited from Middle Chinese, \emph{beanH} and \emph{pjop} (Guangyun, in Baxter's transcription), which glosses the first 'to provide' and the second 'a rule'.
\end{cousinweb}

\begin{grammarlens}[title={What you've built}]
One more word for saying no and asking.
\end{grammarlens}

\subsection*{Your turn}
\begin{itemize}
\raggedright
  \item \emph{Say it:} \emph{bànfǎ}
  \item \emph{Say it:} \emph{bànfǎ}, once more
  \item \emph{Say it:} \emph{wǒ méiyǒu bànfǎ} — There's nothing I can do.
  \item \emph{Recall:} say the Mandarin for to tell someone, then the Mandarin for to grasp a meaning, then say \emph{bànfǎ} again
\end{itemize}

\begin{tcolorbox}[breakable,colback=teal!4,colframe=teal!35!black,title={Before you move on}]
What does bànfǎ mean? (\textquotedblleft{}A way to deal with it\textquotedblright{}.) Say it once more.
\end{tcolorbox}
\section[zěnme bàn]{zěnme bàn — what should I do?}

\label{lesson:ZH-C214-zenmeban}



\begin{quote}
Before the new one: say the Mandarin for to grasp a meaning, then the Mandarin for a way to deal with it.
\end{quote}

\subsection*{You'll want to know: zěnme bàn}
\emph{zěnme bàn} — \textquotedblleft{}what should I do?\textquotedblright{}.

\begin{quote}
\emph{wǒ de hùzhào diū le zěnme bàn} — I've lost my passport. What should I do?
\end{quote}

Its characters come later; for now it is a word for the ear.

\begin{cousinweb}
\emph{zěnme bàn} is built from \emph{zěnme} 'how' and \emph{bàn} 'to handle'. The last is inherited from Middle Chinese \emph{beanH} (Guangyun, in Baxter's transcription); no dictionary we use records an older form of \emph{zěn}, so the history of \emph{zěnme} is not settled.
\end{cousinweb}

\begin{grammarlens}[title={What you've built}]
One more word for saying no and asking.
\end{grammarlens}

\subsection*{Your turn}
\begin{itemize}
\raggedright
  \item \emph{Say it:} \emph{zěnme bàn}
  \item \emph{Say it:} \emph{zěnme bàn}, once more
  \item \emph{Say it:} \emph{wǒ de hùzhào diū le zěnme bàn} — I've lost my passport. What should I do?
  \item \emph{Recall:} say the Mandarin for to grasp a meaning, then the Mandarin for a way to deal with it, then say \emph{zěnme bàn} again
\end{itemize}

\begin{tcolorbox}[breakable,colback=teal!4,colframe=teal!35!black,title={Before you move on}]
What does zěnme bàn mean? (\textquotedblleft{}What should I do?\textquotedblright{}.) Say it once more.
\end{tcolorbox}
